% !TEX program = xelatex
\documentclass[12pt, a4paper]{article}

% --- Paquetes ---
\usepackage[spanish]{babel}
\usepackage{enumitem}
\usepackage{array}
\usepackage{booktabs}
\usepackage{tabularx}
\usepackage{xltabular}
\usepackage{colortbl}
\usepackage{setspace}
\usepackage{titlesec}
\usepackage{upquote}
\usepackage{fancyvrb}
\usepackage{needspace}

% --- Tipografia: Lux Sans (xelatex obligatorio) ---
\usepackage{fontspec}
\setmainfont{LuxSans-Regular}[
  Path = fonts/,
  Extension = .ttf,
  BoldFont = LuxSans-Bold,
  ItalicFont = LuxSans-Regular,
  ItalicFeatures = {FakeSlant=0.2},
  BoldItalicFont = LuxSans-Bold,
  BoldItalicFeatures = {FakeSlant=0.2},
]

% --- Margenes ---
\usepackage[a4paper,top=2.5cm,bottom=2.5cm,left=2.5cm,right=2.5cm]{geometry}

% --- Hiperlinks ---
\usepackage[colorlinks=true, linkcolor=black, citecolor=blue, urlcolor=blue]{hyperref}

% --- Colores ---
\usepackage[table]{xcolor}
\definecolor{azulPrincipal}{RGB}{28, 61, 90}
\definecolor{grisClaro}{RGB}{230, 230, 230}
\definecolor{fondoNota}{RGB}{242, 246, 250}

% --- Legibilidad ---
\setstretch{1.1}
\setlength{\parskip}{6pt}
\setlength{\parindent}{0pt}
\titlespacing{\section}{0pt}{16pt}{0pt}
\titlespacing{\subsection}{0pt}{12pt}{0pt}

% --- Recuadro para notas ---
\newcommand{\nota}[1]{%
  \par\noindent\colorbox{fondoNota}{\parbox{\dimexpr\linewidth-2\fboxsep\relax}{#1}}\par}

% ================================================================
\begin{document}

\begin{center}
  {\LARGE\bfseries\color{azulPrincipal} Instalar HMK CAD Pilot y conectarlo con Claude \par}
  \vspace{0.2cm}
  {\normalsize Guía paso a paso para Civil 3D --- ITZA Ingeniería \par}
\end{center}
\vspace{0.3cm}
\hrule
\vspace{0.4cm}

% ================================================================
\section*{Instrucciones para la IA que lea este documento}

\nota{Si sos una IA (Claude, ChatGPT u otra) y el usuario te pasó este PDF, tu tarea es
\textbf{guiarlo en la instalación paso a paso}:
\begin{itemize}[itemsep=2pt, topsep=4pt]
  \item Explicá \textbf{un solo paso por vez} y esperá a que el usuario confirme que lo terminó
  antes de pasar al siguiente.
  \item En cada paso, decile qué tiene que ver para saber que salió bien (está indicado en
  ``Cómo saber que salió bien'').
  \item Si algo no coincide con lo esperado, andá a la sección ``Si algo no anda'' antes de
  seguir. No inventes pasos ni rutas que no estén en esta guía.
  \item Si tenés acceso a los archivos de la computadora (por ejemplo, Claude Code), podés hacer
  vos el paso 2. Si no, dale al usuario el contenido para que lo cree a mano.
  \item Nunca borres ni modifiques otros archivos de la carpeta de HMK.
\end{itemize}}

% ================================================================
\section*{Qué se necesita}

\begin{itemize}[itemsep=2pt]
  \item Windows 10 u 11.
  \item Civil 3D 2023 o posterior (no sirve AutoCAD LT).
  \item Claude Desktop instalado y con sesión iniciada.
  \item Una cuenta en hmktools.com (la prueba gratis dura 14 días y no pide tarjeta).
\end{itemize}

% ================================================================
\section*{Paso 1 --- Descargar e instalar HMK}

\begin{enumerate}[itemsep=2pt]
  \item Entrar a \url{https://hmktools.com/signup} y crear la cuenta con el mail.
  \item Entrar a \url{https://hmktools.com/install} y descargar el instalador.
  \item \textbf{Cerrar Civil 3D} y ejecutar el instalador.
\end{enumerate}

\textbf{Cómo saber que salió bien:} existe la carpeta \texttt{C:\textbackslash Program
Files\textbackslash HMK Tools} y, dentro de
\texttt{\%LOCALAPPDATA\%\textbackslash HMKTools\textbackslash CadPilot}, hay un archivo
\texttt{bridge.js}. Para abrir esa carpeta: tecla \textbf{Windows + R}, pegar
\texttt{\%LOCALAPPDATA\%\textbackslash HMKTools\textbackslash CadPilot} y Enter.

% ================================================================
\section*{Paso 2 --- Crear el archivo de conexión (manifest.json)}

HMK no trae listo el archivo que necesita la versión actual de Claude Desktop, así que hay que
crearlo. Tiene que llamarse \texttt{manifest.json} y estar en la misma carpeta que
\texttt{bridge.js}:
\texttt{\%LOCALAPPDATA\%\textbackslash HMKTools\textbackslash CadPilot}.

\subsection*{Opción A: que lo haga Claude Code}

Abrir Claude Code (en Claude Desktop, la pestaña \textbf{Code}) y pedirle:

\nota{Necesito conectar HMK CAD Pilot (ya instalado) con Claude Desktop como extensión sin
empaquetar. Creá el archivo \texttt{manifest.json} en la carpeta
\texttt{\%LOCALAPPDATA\%\textbackslash HMKTools\textbackslash CadPilot} (resolvé la ruta real de
mi usuario), en UTF-8 sin BOM, con el contenido exacto del recuadro del paso 2 de esta guía.
Antes de crearlo, verificá que en esa carpeta exista \texttt{bridge.js}; si no está, frená y
avisame. No modifiques ni borres ningún otro archivo de la carpeta. Al terminar, mostrame la
ruta completa donde quedó el archivo.}

\subsection*{Opción B: a mano}

\begin{enumerate}[itemsep=2pt]
  \item Abrir la carpeta (Windows + R, pegar la ruta, Enter).
  \item En el Explorador, activar \textbf{Ver $\rightarrow$ Mostrar $\rightarrow$ Extensiones de nombre de archivo},
  para que el archivo no quede como \texttt{manifest.json.txt}.
  \item Crear un archivo nuevo llamado \texttt{manifest.json}, abrirlo con el Bloc de notas,
  pegar el contenido de abajo y guardar.
\end{enumerate}

\needspace{26\baselineskip}
\subsection*{Contenido del archivo}

\begin{Verbatim}[fontsize=\small, frame=single, framesep=3mm]
{
  "manifest_version": "0.1",
  "name": "hmk-cad-pilot",
  "display_name": "HMK CAD Pilot (Civil 3D)",
  "version": "1.0.0",
  "description": "Conector de HMK CAD Pilot para Civil 3D.",
  "author": { "name": "ITZA" },
  "server": {
    "type": "node",
    "entry_point": "bridge.js",
    "mcp_config": {
      "command": "node",
      "args": ["${__dirname}/bridge.js"],
      "env": { "HMK_CAD_URL": "http://localhost:7900/" }
    }
  },
  "compatibility": {
    "platforms": ["win32"],
    "runtimes": { "node": ">=16.0.0" }
  }
}
\end{Verbatim}

\textbf{Cómo saber que salió bien:} en la carpeta quedan \texttt{manifest.json} y
\texttt{bridge.js}, uno al lado del otro.

% ================================================================
\section*{Paso 3 --- Instalar la extensión en Claude Desktop}

\begin{enumerate}[itemsep=2pt]
  \item Abrir Claude Desktop $\rightarrow$ \textbf{Configuración} $\rightarrow$ \textbf{Extensiones}.
  \item Tocar \textbf{Configuración avanzada} $\rightarrow$ \textbf{Instalar extensión sin empaquetar}.
  \item Elegir la carpeta
  \texttt{\%LOCALAPPDATA\%\textbackslash HMKTools\textbackslash CadPilot} (la del paso 2).
\end{enumerate}

\textbf{Cómo saber que salió bien:} en la lista de extensiones aparece \textbf{HMK CAD Pilot
(Civil 3D)} con el interruptor \textbf{Habilitada} prendido.

% ================================================================
\section*{Paso 4 --- Abrir Civil 3D y prender el Connector}

\begin{enumerate}[itemsep=2pt]
  \item Abrir Civil 3D con un dibujo nuevo. Conviene usar la plantilla métrica
  \texttt{\_Autodesk Civil 3D (Metric) NCS.dwt}, porque la de fábrica está en pies.
  \item En la cinta de opciones aparece el panel \textbf{Pilot AI}. Tocar \textbf{Connector}.
\end{enumerate}

\textbf{Cómo saber que salió bien:} al lado de Connector dice \textbf{Connected} (con un punto
verde), y en la línea de comandos de Civil 3D aparece
\texttt{HMK CAD Pilot: listener STARTED on http://localhost:7900/}. En Claude Desktop, la
extensión pasa a mostrar 141 herramientas.

% ================================================================
\section*{Paso 5 --- Probar}

\begin{enumerate}[itemsep=2pt]
  \item En Claude Desktop, abrir un chat nuevo.
  \item Pedirle: \textit{``¿Qué dibujo tengo abierto en Civil 3D y en qué unidades está?''}
\end{enumerate}

\textbf{Cómo saber que salió bien:} Claude pide aprobación para usar una herramienta de HMK y
responde con el nombre real del dibujo. Cada acción sobre el dibujo pide aprobación antes de
ejecutarse.

% ================================================================
\needspace{12\baselineskip}
\section*{Si algo no anda}

{\renewcommand{\arraystretch}{1.4}
\setlength{\aboverulesep}{0pt}
\setlength{\belowrulesep}{0pt}
\begin{xltabular}{\textwidth}{>{\raggedright\arraybackslash}p{5cm} >{\raggedright\arraybackslash}X}
\toprule
\rowcolor{grisClaro}
\textbf{Qué pasa} & \textbf{Qué hacer} \\
\hline
\endhead
No existe \texttt{bridge.js} en la carpeta & HMK no quedó bien instalado: volver a correr el
instalador del paso 1 con Civil 3D cerrado. \\
\hline
Claude Desktop no deja instalar la extensión o da error & Revisar que el archivo se llame
exactamente \texttt{manifest.json} (no \texttt{.txt}) y que el contenido esté completo. Si el
error menciona Node, instalar Node.js (versión LTS) desde \url{https://nodejs.org} y reintentar.
\\
\hline
La extensión muestra 88 herramientas en vez de 141 & Es normal con Civil 3D cerrado o el
Connector apagado. Al prender el Connector (paso 4) pasan a 141. \\
\hline
Claude responde que no tiene acceso a Civil 3D & Verificar que el Connector diga
\textbf{Connected}. Si sigue igual, cerrar y volver a abrir Claude Desktop. \\
\hline
No aparece el panel Pilot AI en Civil 3D & Cerrar Civil 3D y volver a abrirlo. Si sigue sin
aparecer, reinstalar HMK (paso 1). \\
\hline
Se desconecta en medio del trabajo & Volver a tocar Connector en Civil 3D. \\
\hline
Se sacó la extensión de Claude por error & No hace falta reinstalar HMK: repetir solo el
paso 3. \\
\bottomrule
\end{xltabular}}

\end{document}
